\honest{Is Humanism A Religion-Substitute?}{Eliezer Yudkowsky}{2008}

Long before the Wright Brothers, people dreamed of flying with magic potions. The wish to fly was fine; only the potions were irrational.

Suppose a brain scan showed that my brain, watching a space shuttle launch (wanting to visit space is not ``realistic,'' but it is a lawful dream), looked like a devout Christian's watching a nativity scene. Maybe it would, maybe not. Christians and some atheists would gloat: ``space travel is your religion!'' That draws the wrong category boundary. Some people once tried to fly by magic, and that does not taint looking down on clouds from a plane. To assume that every nonbeliever has a hole in the mind that wants filling is what I call ``theomorphism.''

To be fair, some atheists do have religion-shaped holes. I have seen their substitutes, and they are ``invariably awful.'' I call them ``hymns to the nonexistence of God'': ``Hail, oh unintelligent universe,'' and so on. They will ``without exception, suck.'' I name none. They are bad because they are imitative: their writers have no motive but ``a vague feeling that since churches have hymns, they ought to have one too.'' Genuine religious art, by contrast, ``is not an imitation of anything'': its writers often felt strongly, wrote honestly and worked hard at prosody and imagery, which gave it artistic integrity. \nb{The hymn tradition imitates too. Isaac Watts, credited by many with introducing hymns to English churches, titled a 1719 psalter \textsc{Psalms of David: Imitated in the Language of the New Testament}. So the post's case rests on the motives it assigns to the unnamed writers.}

My test for post-religious practice: would it make sense if religion had never existed? Hymns to God's nonexistence fail. Weddings pass, as long as no one lectures on God: in a world without religion, a wedding would be about love, children, commitment and devotion, and no one would mention God. \nb{Humanist weddings had been legal in Scotland since 2005; by 2017 the Humanist Society Scotland conducted more weddings than the Church of Scotland.} Tears at a shuttle launch pass too, so even a scan showing ``religion''-associated areas lighting up would not make them a substitute. A good atheistic hymn is a song about anything worth singing about that is not religious. Nor should you avoid things just because they resemble religion, since ``reversed stupidity is not intelligence.'' The world's greatest idiot may say the Sun is shining, but that does not make it dark. Believe accurately, then feel accordingly: if watching a rocket rise makes you want to sing, ``write the song, dammit.'' Tears at a launch mean that my caring is bound into the real world.

If God spoke plainly and answered prayers reliably, God would become ``one more boringly real thing,'' no more worth believing in than the postman. A real God would destroy ``the inner uncertainty that brings forth outward fervor in compensation.'' If everyone believed, being ``one of the elect'' would stop being special. I give no evidence for either claim. \nb{Experiments by McGregor and colleagues (2001) found that people made uncertain responded with more extreme conviction; the abstract concerns social issues and group bias, not religion.} Awe at a real shuttle has no such weakness: I can see it rise without losing the awe, and everyone else believing in shuttles does not make them less special.

``The choice between God and humanity is not just a choice of drugs. Above all, humanity actually exists.'' \nb{The word ``humanism'' appears only in the title. Organized humanism's first manifesto, in 1933, presented humanism as a new religion. The post does not discuss it.}
